\documentclass[12pt]{article}

\usepackage{tikz} % картинки в tikz
\usepackage{microtype} % свешивание пунктуации

\usepackage{array} % для столбцов фиксированной ширины

\usepackage{indentfirst} % отступ в первом параграфе

\usepackage{sectsty} % для центрирования названий частей
\allsectionsfont{\centering}

\usepackage{amsmath} % куча стандартных математических плюшек

\usepackage{comment}

\usepackage[top=2cm, left=1.2cm, right=1.2cm, bottom=2cm]{geometry} % размер текста на странице

\usepackage{lastpage} % чтобы узнать номер последней страницы

\usepackage{enumitem} % дополнительные плюшки для списков
%  например \begin{enumerate}[resume] позволяет продолжить нумерацию в новом списке
\usepackage{caption}


\usepackage{fancyhdr} % весёлые колонтитулы
\pagestyle{fancy}
\lhead{Работа над ошибками :)}
\chead{}
\rhead{Лекция ПСМО 2020-09-25}
\lfoot{}
\cfoot{}
\rfoot{\thepage/\pageref{LastPage}}
\renewcommand{\headrulewidth}{0.4pt}
\renewcommand{\footrulewidth}{0.4pt}



\usepackage{todonotes} % для вставки в документ заметок о том, что осталось сделать
% \todo{Здесь надо коэффициенты исправить}
% \missingfigure{Здесь будет Последний день Помпеи}
% \listoftodos - печатает все поставленные \todo'шки


% более красивые таблицы
\usepackage{booktabs}
% заповеди из докупентации:
% 1. Не используйте вертикальные линни
% 2. Не используйте двойные линии
% 3. Единицы измерения - в шапку таблицы
% 4. Не сокращайте .1 вместо 0.1
% 5. Повторяющееся значение повторяйте, а не говорите "то же"



\usepackage{fontspec}
\usepackage{polyglossia}

\setmainlanguage{russian}
\setotherlanguages{english}

% download "Linux Libertine" fonts:
% http://www.linuxlibertine.org/index.php?id=91&L=1
\setmainfont{Linux Libertine O} % or Helvetica, Arial, Cambria
% why do we need \newfontfamily:
% http://tex.stackexchange.com/questions/91507/
\newfontfamily{\cyrillicfonttt}{Linux Libertine O}

\AddEnumerateCounter{\asbuk}{\russian@alph}{щ} % для списков с русскими буквами
\setlist[enumerate, 2]{label=\asbuk*),ref=\asbuk*}

%% эконометрические сокращения
\DeclareMathOperator{\Cov}{\mathbb{C}ov}
\DeclareMathOperator{\Corr}{\mathbb{C}orr}
\DeclareMathOperator{\Var}{\mathbb{V}ar}
\DeclareMathOperator{\E}{\mathbb{E}}
\DeclareMathOperator{\tr}{trace}
\def \hb{\hat{\beta}}
\def \hs{\hat{\sigma}}
\def \htheta{\hat{\theta}}
\def \s{\sigma}
\def \hy{\hat{y}}
\def \hY{\hat{Y}}
\def \v1{\vec{1}}
\def \e{\varepsilon}
\def \he{\hat{\e}}
\def \z{z}
\def \hVar{\widehat{\Var}}
\def \hCorr{\widehat{\Corr}}
\def \hCov{\widehat{\Cov}}
\def \cN{\mathcal{N}}


\begin{document}
 
\includegraphics[scale=0.8]{CaptainFacepalm.png}

В конце лекции я написал формулу:

\begin{multline}
 Q(\theta, \theta_{old}) = 
 \sum_{i=1}^n P(z_i = 1\mid x, \theta_{old}) \left( -\frac{1}{2}\ln(2\pi\sigma_1^2) - \frac{1}{2}\frac{(x_i - \mu_1)^2}{\sigma_1^2} \right) + \\ 
 (1- P(z_i = 1\mid x, \theta_{old})) \left( -\frac{1}{2}\ln(2\pi\sigma_2^2) - \frac{1}{2}\frac{(x_i - \mu_2)^2}{\sigma_2^2} \right)
\end{multline}


Это неверно! Снимаем лапшу с ушей!

Так как нам нужен логарифм от совместной плотности в ожидании $E(\ln p(x,z|\theta)\mid x, \theta_{old})$, 
то нужно добавить $\ln P(z_i = 1|\theta)$ и $\ln (1 - P(z_i = 1|\theta))$. Плотность домножается, поэтому логарифм прибавляем.

\begin{multline}
 Q(\theta, \theta_{old}) = 
 \sum_{i=1}^n P(z_i = 1\mid x, \theta_{old}) \left( -\frac{1}{2}\ln(2\pi\sigma_1^2) - \frac{1}{2}\frac{(x_i - \mu_1)^2}{\sigma_1^2}  \color{red}+ \ln p_1 \color{black}\right) + \\ 
 (1- P(z_i = 1\mid x, \theta_{old})) \left( -\frac{1}{2}\ln(2\pi\sigma_2^2) - \frac{1}{2}\frac{(x_i - \mu_2)^2}{\sigma_2^2} \color{red}+ \ln (1- p_1) \color{black}\right)
\end{multline}
   

В результате оптимальные $\mu_1^{new}$, $\mu_2^{new}$, $\sigma^2_{1,new}$, $\sigma^2_{2,new}$ не меняются! 
Это меня и смутило, я ожидал, что при ошибке в формуле поменяются оптимумы.

А вот оптимальное $p_1^{new}$ оказывается результатом честной оптимизации, а не эвристики, как я сказал.


\end{document}
